
%%%%%%%%%%%%%%%%%%%%%%%%%%
\documentclass[12pt, a4paper]{article}
\input{systeme}
\usepackage[a4paper, top=4cm,left=2.5cm, right=2.3cm, bottom=2.7cm]{geometry}
\usepackage{amsmath, amssymb, array, multicol, enumitem}
\usepackage{pdfpages}
\usepackage{comment, graphicx}  
\usepackage{xparse}
\DeclareGraphicsExtensions{.png,.jpg}
\usepackage{pgffor}  % for loops in latex
\usepackage{ifthen}  % conditional latex compilation
\usepackage{todonotes}  % notes in bright orange
\usepackage[colorlinks=true]{hyperref}  % hyperlinks
\AtBeginDocument{% coloring of links
	\hypersetup{pdfborder={0 0 0},
		urlcolor=blue}%
}
\usepackage{fontspec}  % Hebrew fonts
\usepackage{xunicode}  % Hebrew character encoding
\usepackage{polyglossia}  % Hebrew Latex
\usepackage[absolute]{textpos}   % Proper positioning of images, e.g. BGU logo

\setdefaultlanguage{hebrew}
\setotherlanguage{english}

\newfontfamily\hebrewfont[Script=Hebrew]{Times New Roman}

\makeatletter
\def\tagform@#1{\maketag@@@{(\ignorespaces#1\unskip)}}
\makeatother


\makeatletter
\renewcommand{\SepMark}[1]{\def\@SepMark{#1}}
\makeatother
\SepMark{.}
\robustify\hebrewnumeral
\robustify\Hebrewnumeral
\robustify\Hebrewnumeralfinal

\def\urladdrname{{\itshape כתובת רשת}}
\def\emailaddrname{{\itshape כתובת דואל}}

\makeatletter
\let\orig@MakeUppercase \MakeUppercase
\makeatother
\providecommand{\MakeUppercase}[1]{#1}

\fontsize{14}{15}  

\setdefaultlanguage{hebrew}
\setotherlanguage{english}
\thispagestyle{empty}
\selectlanguage{hebrew}

\newcommand{\cc}{\mathbb{C}}
\newcommand{\CC}{\mathbb{C}}
\newcommand{\rr}{\mathbb{R}}
\newcommand{\RR}{\mathbb{R}}
\newcommand{\nn}{\mathbb{N}}
\newcommand{\zz}{\mathbb{Z}}
\newcommand{\ZZ}{\mathbb{Z}}
\newcommand{\qq}{\mathbb{Q}}
\newcommand{\FF}{\mathbb{F}}
\newcommand{\QQ}{\mathbb{Q}}
\DeclareMathOperator{\vol}{vol}
\newcommand{\abs}[1]{\left|#1\right|}

\DeclareMathOperator{\Adj}{Adj}
\DeclareMathOperator{\deter}{det}
\DeclareMathOperator{\Img}{Im}
\DeclareMathOperator{\Ker}{Ker}
\DeclareMathOperator{\rank}{rank}
\DeclareMathOperator{\Span}{Span}
\DeclareMathOperator{\tr}{tr}


\newcommand{\calB}{\mathcal{B}}
\newcommand{\calC}{\mathcal{C}}
\newcommand{\cD}{\mathcal{D}}
\newcommand{\calE}{\mathcal{E}}


%%%% new for complex functions

\newcommand{\Arg}{\mathrm{Arg}}
\newcommand{\ol}{\overline}
\newcommand{\wt}{\widetilde}
\newcommand{\wh}{\widehat}
\newcommand{\re}{\mathrm{Re}}
\newcommand{\im}{\mathrm{Im}}
\newcommand{\goes}{\longrightarrow}
\newcommand{\cO}{\mathcal{O}}
\newcommand{\Log}{\mathrm{Log}}


%%%%%%%%%%%%%%%%%%%%%%%%%%%%%%%%%%%%


\begin{document}
\selectlanguage{hebrew}

\newcommand{\semester}{%
	סמסטר א' תשפ״ה
}
\newcommand{\course}{%
הסתברות לסטטיסטיקה}
\newcommand{\coursenumber}{9321}
\newcommand{\doctype}{%
\boxed{\text{
 דף  תרגילים 1 (להגשה עד 03.12 בשעה 23:59)}}
}
\newcommand{\hwknumber}{}
\newcommand{\chaptype}{%
	פרק }
\newcommand{\chapnum}{\textenglish{I}}

\newcommand{\chapname}{%
אאא } 
	
\thispagestyle{empty}
\begin{textblock*}{9cm}[0.5,0.5](90mm,3cm)
	\includegraphics[width=2in]{logo}
\end{textblock*}	

\begin{center}
	\large	
המחלקה למתמטיקה - 
	\semester  \\
	{\large \bf \course - 201.1.9921}\\
	\medskip
	{\Large  \textbf{\doctype \hwknumber}  }\\
\end{center}
\vspace{0.3cm}

%%%%%%%%%%%%%%%%%%%%%%%%%%%%%%%%%%%%

\noindent\rule{17cm}{0.4pt}

\begin{enumerate}


\item
גפן הולכת למכולת השכונתית ומתלבטת איזה מבין שלושת הסוגים הבאים של חלב לקנות: סויה, שקדים או שיבולת-שועל.
נסמן את המאורעות הבאים:

$A$
-
גפן קנתה חלב סויה,

$B$
-
גפן קנתה חלב שקדים,

$C$
גפן קנתה חלב שיבולת-שועל.

תארו כל אחד מהמאורעות (קבוצות) הבאים בעזרת
$A,B,C$
והפעולות חיתוך, חיבור, הפרש ומשלים של קבוצות. בנוסף ציירו את דיאגרמת ון המתאימה בכל אחד מהסעיפים:
\begin{enumerate}
\item
גפן קונה את שלושת הסוגים.
\item
גפן קונה חלב סויה  אבל לא קונה שקדים ולא קונה שיבולת-שועל.
\item
גפן קונה חלב שקדים, ועוד בדיוק אחד מבין השניים האחרים.
\item
גפן קונה לכל היותר סוג חלב אחד.
\item
גפן קונה בדיוק סוג חלב אחד.
\item
גפן קונה בדיוק שני סוגי חלב.
\item
גפן לא קונה חלב כלל.
\item 
גפן קונה לפחות שני סוגים שונים של חלב.
\end{enumerate}

\noindent\rule{15cm}{0.4pt}

\item
פרופסור אקס רוצה לסדר במשרדו החדש
$10$
ספרים 
\underline{שונים}
על המדף:

$4$
ספרי היסטוריה,
$3$
ספרי מדע בדיוני,
$2$
ספרי רפואה,
וספר סטטיסטיקה אחד.
\begin{enumerate}
\item
בכמה דרכים הוא יכול לעשות זאת?
\item
בכמה דרכים יוכל לעשות זאת, אם הוא רוצה שספר הסטטיסטיקה יהיה באחד מהקצוות של המדף?
\item
אם הוא רוצה לסדר את הספרים כך שכל שהספרים באותו הנושא יהיו צמודים זה לזה, בכמה דרכים יוכל לעשות זאת?
\item
אחיינית של פרופסור אקס מגיעה לביקור, וכמתנה ליום הולדתה השביעי הוא רוצה להביא לה במתנה
שלושה
ספרים מתוך העשרה שיש לו. כמה מתנות שונות יכול הפרופסור לתת לאחייניתו האהובה?
\item
בכל יום ראשון בוחר הפרופסור שני ספרים ושם אותם בתיבת הדואר של התאגדות הסטודנטים. 
מהיותו פרופסור מסודר, הוא מנהל רשימה בה רושם את שמות שני הספרים שבחר בכל שבוע.
\begin{enumerate}
\item[(1)]
בהנחה כי הסטודנטים לא  מחזירים לפרופסור את הספרים,
כמה רשימות שונות יכול לנהל הפרופסור עד שיגמרו לו הספרים על המדף? (כלומר, בתום 
$5$
שבועות)
\item[(2)]
בהנחה
שהסטודנטים מחזירים לפרופסור את הספרים עד סוף השבוע, בכל שבוע, כמה רשימות שונות יכולות היו להיות לפרופסור בתום
$5$
שבועות?
\end{enumerate}
\end{enumerate}

\item
 בשאלה זו נעסוק במקדם הבינומי
 $$\binom{n}{k}=\frac{n!}{k!(n-k)!}$$
 כאשר משמעותו הקומבינטורית היא ״מספר הדרכים בהם ניתן לבחור
 $k$
 עצמים מתוך
 $n$
 עצמים שונים״.
\begin{enumerate}
\item
חשבו, ללא מחשבון:
\begin{enumerate}
\begin{multicols}{3}
\item[(1)]
$\binom{4}{2}$
\item[(2)]
$\binom{100}{1}$
\item[(3)]
$\binom{20}{2}$
\end{multicols}
\end{enumerate}
\item
הראו כי
לכל
$1\le k\le n$
 מתקיימת הזהות הבאה:
$$.k\cdot \binom{n}{k}
=n\cdot \binom{n-1}{k-1}$$
\item
נסמן
$\Omega=\{A,B,C\}$.
\begin{enumerate}
\item[(1)]
כמה תתי-קבוצות שונות 
יש ל-$\Omega$?
כתבו את כולן.
\item[(2)]
כמה תתי-קבוצות שונות יש ל-$\Omega$
שהן בגודל זוגי? ($0$
נחשב זוגי)
\item[(3)]
כמה תתי-קבוצות שונות יש ל-$\Omega$
שהן בגודל אי-זוגי? \\
\end{enumerate}

\item
נסמן
$\Omega=\{A,B,C,D,E,F,G,H\}$.
\begin{enumerate}
\item[(1)]
כמה תתי-קבוצות שונות בגודל
$2$
יש ל-$\Omega$?
\item[(2)]
כמה תתי-קבוצות שונות יש ל-$\Omega$?
\item[(3)]
כמה תתי-קבוצות שונות יש ל-$\Omega$
שהן בגודל זוגי? ($0$
נחשב זוגי)
\item[(4)]
כמה תתי-קבוצות שונות יש ל-$\Omega$
שהן בגודל אי-זוגי?\\ 
\end{enumerate}
\item
היעזרו בנוסחת הבינום של ניוטון שנלמדה בכיתה
$$(a+b)^n=
\sum_{k=0}^{n}
\binom{n}{k}a^kb^{n-k}$$
והראו כי לכל קבוצה סופית
$\Omega$
מספר תתי-הקבוצות השונות
של
$\Omega$
שהן בגודל זוגי תמיד יהיה שווה למספר תתי-הקבוצות השונות 
של
$\Omega$
שהן בגודל אי-זוגי.
(רמז: סמנו
$|\Omega|=n$
ובחרו את
$a,b$ 
להיות....)
\end{enumerate}

\noindent\rule{15cm}{0.4pt}

\item
\begin{enumerate}
\item
נסו לתת הסבר קומבינטורי לשיוויון הבא:
$$\binom{30}{13}
=\sum_{k=0}^{13}
\binom{10}{k}\cdot \binom{20}{13-k}$$
(רמז:
חשבו על
גפן שצריכה לארוז 
$13$
זוגות מכנסיים 
ביציאתה לחופשה, כאשר בארון יש לה
$30$
זוגות מכנסיים, מתוכם 
$10$
ארוכים
ו-$20$
קצרים. בכמה דרכים תוכל לעשות זאת?)
\item
נסו לחשוב על שיקול דומה לסעיף הקודם, ולהוכיח כי לכל
$n\in\nn$
מתקיימת הזהות הבאה:
$$\binom{2n}{n}=
\sum_{k=0}^{n}
\binom{n}{k}^2$$
(הכוונה: כמה זוגות מכנסיים יש לגפן בארון כעת, וכמה היא צריכה לארוז?)
\end{enumerate}
\end{enumerate}


\end{document}





